\documentclass[11pt,letterpaper]{article}
\usepackage[utf8]{inputenc}
\usepackage[T1]{fontenc}
\usepackage[spanish,es-nodecimaldot]{babel}
\usepackage{lmodern,microtype}
\usepackage{amsmath,amssymb}
\usepackage[margin=2.2cm]{geometry}
\usepackage{booktabs,array,enumitem}
\usepackage[dvipsnames]{xcolor}
\usepackage{tcolorbox}
\usepackage{pgfplots}
\pgfplotsset{compat=1.17}
\usepgfplotslibrary{fillbetween}
\usepackage{hyperref}
\hypersetup{colorlinks=true,linkcolor=NavyBlue,urlcolor=NavyBlue}

\tcbuselibrary{skins,breakable}
\newtcolorbox{formula}[1]{colback=NavyBlue!4,colframe=NavyBlue!70,fonttitle=\bfseries,title={#1},breakable}
\newtcolorbox{nota}[1]{colback=OliveGreen!5,colframe=OliveGreen!70,fonttitle=\bfseries,title={#1},breakable}
\newtcolorbox{paso}[1]{colback=white,colframe=gray!60,fonttitle=\bfseries\small,coltitle=black,colbacktitle=gray!15,title={#1},breakable}

\setlength{\parindent}{0pt}
\setlength{\parskip}{0.45em}
\renewcommand{\arraystretch}{1.15}
\newcommand{\SSR}{\mathrm{SSR}}
\newcommand{\SSE}{\mathrm{SSE}}
\newcommand{\SST}{\mathrm{SST}}

\title{\textbf{Las pruebas $t$ y $F$ paso a paso}\\[4pt]
\large Tres ejercicios resueltos con datos sintéticos\\[2pt]
\normalsize Fundamentos de Econometría -- Universidad Panamericana -- Otoño 2026}
\author{Dr. César Galindo}
\date{}

\begin{document}
\maketitle

\begin{nota}{Objetivo de la guía}
Resolver dos pruebas $F$ en regresión lineal múltiple y una prueba $t$ individual (tema 4.3 de
las notas del segundo parcial). Se presentan las propiedades y usos de ambas distribuciones, los
cálculos de los estadísticos, las reglas de decisión y la relación entre las pruebas. Los cálculos
se muestran con fórmulas y números intermedios para que puedan reproducirse con calculadora.
\end{nota}

\section{Recetario}

\subsection{Notación}

El modelo poblacional con $k$ regresoras y $n$ observaciones es
\[
Y_i=\beta_0+\beta_1X_{1i}+\cdots+\beta_kX_{ki}+\varepsilon_i,\qquad i=1,\dots,n,
\qquad\text{o bien}\qquad \mathbf y=X\boldsymbol\beta+\boldsymbol\varepsilon .
\]
El estimador de MCO es $\hat{\boldsymbol\beta}=(X'X)^{-1}X'\mathbf y$, los valores ajustados son
$\hat Y_i$ y los residuos $\hat e_i=Y_i-\hat Y_i$. Se usan tres sumas de cuadrados:
\[
\SST=\sum_{i=1}^n (Y_i-\bar Y)^2,\qquad
\SSE=\sum_{i=1}^n (\hat Y_i-\bar Y)^2,\qquad
\SSR=\sum_{i=1}^n \hat e_i^{\,2},\qquad \SST=\SSE+\SSR .
\]
\textbf{Convención:} $\SSR$ es la suma de cuadrados de los \emph{residuos} (lo que el modelo no
explica) y $\SSE$ la suma de cuadrados \emph{explicada}, como en las notas. Algunos libros usan las
letras al revés; lo que importa es que en la prueba $F$ entra siempre la suma de cuadrados
\emph{residual}. Además $R^2=1-\SSR/\SST=\SSE/\SST$.

\subsection{Modelo restringido y no restringido}

\begin{itemize}[leftmargin=1.3em]
\item \textbf{No restringido (NR):} se estiman libremente los $k+1$ parámetros. Su suma de cuadrados
residual es $\SSR_{NR}$.
\item \textbf{Restringido (R):} se imponen las $q$ restricciones de $H_0$ (por ejemplo
$\beta_2=\beta_3=0$, lo que equivale a quitar $X_2$ y $X_3$) y se vuelve a estimar. Su suma de
cuadrados residual es $\SSR_R$.
\end{itemize}
Como el modelo restringido tiene menos libertad para ajustar los datos, siempre
$\SSR_R\ge \SSR_{NR}$. La pregunta de la prueba es si el aumento $\SSR_R-\SSR_{NR}$ es grande en
comparación con el ruido residual del modelo no restringido.

\begin{formula}{Estadístico $F$}
\[
F=\frac{(\SSR_R-\SSR_{NR})/q}{\SSR_{NR}/(n-k-1)}
\;\overset{H_0}{\sim}\;F_{q,\;n-k-1},
\]
donde $q$ es el número de restricciones, $n$ el tamaño de muestra y $k$ el número de regresoras
del modelo \emph{no restringido} (sin contar la constante). Si ambos modelos tienen la misma
variable dependiente, la forma equivalente con $R^2$ es
\[
F=\frac{(R^2_{NR}-R^2_{R})/q}{(1-R^2_{NR})/(n-k-1)} .
\]
\textbf{Caso particular (prueba global):} si $H_0:\beta_1=\cdots=\beta_k=0$, el modelo restringido
es $Y_i=\beta_0+\varepsilon_i$, su residual es $\SSR_R=\SST$ y $R^2_R=0$:
\[
F=\frac{\SSE/k}{\SSR/(n-k-1)}=\frac{R^2/k}{(1-R^2)/(n-k-1)} .
\]
\end{formula}

\subsection{Distribuciones $t$ de Student y $F$: propiedades y uso}

\subsubsection*{Distribución $t$ de Student}
Si $Z\sim N(0,1)$ y $U\sim\chi^2_{\nu}$ son independientes, entonces
\[
T=\frac{Z}{\sqrt{U/\nu}}\sim t_{\nu}.
\]
La distribución $t_{\nu}$ tiene soporte en toda la recta real y es simétrica alrededor de cero.
Su forma es parecida a la normal estándar, pero sus colas son más pesadas: con pocos grados de
libertad asigna más probabilidad a valores extremos. Al aumentar $\nu$, se aproxima a $N(0,1)$.
Su media es cero si $\nu>1$ y no existe si $\nu\le1$; su varianza es $\nu/(\nu-2)$ si $\nu>2$
y no existe si $\nu\le2$.

\textbf{¿Por qué se usa en regresión?} Para probar una hipótesis individual como
$H_0:\beta_j=\beta_{j,0}$ cuando la varianza de los errores es desconocida. Se estima esa varianza
con los residuos y se divide la diferencia entre el estimador y el valor hipotético por su error
estándar:
\[
t=\frac{\hat\beta_j-\beta_{j,0}}{\operatorname{se}(\hat\beta_j)}.
\]
Bajo los supuestos clásicos de errores normales, independientes y homocedásticos, bajo $H_0$ esta
estadística sigue una $t$ con $n-k-1$ grados de libertad. Para una alternativa bilateral se compara
$|t|$ con el valor crítico $t_{\alpha/2;\,n-k-1}$; valores grandes en cualquiera de las dos colas
son evidencia contra $H_0$. La misma distribución se usa para construir intervalos de confianza
individuales para los coeficientes.

\subsubsection*{Distribución $F$}
Si $U_1\sim\chi^2_{\nu_1}$ y $U_2\sim\chi^2_{\nu_2}$ son independientes, entonces
\[
F=\frac{U_1/\nu_1}{U_2/\nu_2}\sim F_{\nu_1,\nu_2}.
\]
La distribución $F$ sólo toma valores no negativos, suele ser asimétrica hacia la derecha y queda
determinada por sus dos grados de libertad: $\nu_1$ (numerador) y $\nu_2$ (denominador). Se usa en
la prueba de este apunte porque compara el aumento promedio en la suma de cuadrados residual al
imponer las restricciones con la varianza residual estimada en el modelo no restringido. Bajo
$H_0$ y los mismos supuestos clásicos, esa razón sigue una $F_{q,n-k-1}$. Por ello, la región de
rechazo está en la cola derecha: un valor $F$ grande indica que las restricciones empeoran el
ajuste más de lo que cabría atribuir al ruido residual.

\begin{nota}{Cómo elegir entre $t$ y $F$}
La prueba $t$ evalúa una restricción sobre un coeficiente; la prueba $F$ evalúa una o varias
restricciones lineales de manera conjunta. Si se prueba una sola restricción bilateral en el mismo
modelo, ambas pruebas son equivalentes:
\[
F_{1,\nu}=t_{\nu}^{\,2},\qquad
F_{\text{observado}}=t_{\text{observado}}^{\,2}.
\]
El cuadrado elimina el signo de $t$, por eso la equivalencia corresponde a la prueba $t$ bilateral,
no a una prueba unilateral. Con varias restricciones no basta con mirar las pruebas $t$ por
separado: la prueba $F$ también toma en cuenta cómo se estiman conjuntamente los coeficientes.
\end{nota}

\subsection{Los seis pasos de toda prueba $F$}

\begin{enumerate}[leftmargin=1.5em,itemsep=1pt]
\item \textbf{Hipótesis.} Escribir $H_0$ (las $q$ restricciones) y $H_1$ (al menos una no se cumple).
\item \textbf{Modelo no restringido.} Estimar por MCO y calcular $\SSR_{NR}$.
\item \textbf{Modelo restringido.} Imponer $H_0$, estimar y calcular $\SSR_R$.
\item \textbf{Estadístico.} Sustituir en la fórmula de $F$.
\item \textbf{Distribución de referencia y regla de decisión.} Bajo $H_0$ y los supuestos clásicos
(incluida normalidad de los errores), $F\sim F_{q,n-k-1}$. Se rechaza $H_0$ al nivel $\alpha$ si
$F>F_{\alpha;\,q,\,n-k-1}$. La prueba es de una sola cola (derecha): sólo valores grandes de $F$
contradicen $H_0$.
\item \textbf{Conclusión en contexto}, en el lenguaje del problema.
\end{enumerate}

\begin{nota}{Valores críticos que se usan en esta guía}
\centering
\begin{tabular}{lccc}
\toprule
Distribución & $\alpha=0.10$ & $\alpha=0.05$ & $\alpha=0.01$\\
\midrule
$F_{2,7}$ & 3.257 & 4.737 & 9.547\\
$F_{1,7}$ & 3.589 & 5.591 & 12.246\\
$F_{2,8}$ & 3.113 & 4.459 & 8.649\\
$F_{3,8}$ & 2.924 & 4.066 & 7.591\\
$t_7$ (bilateral) & 1.895 & 2.365 & 3.499\\
\bottomrule
\end{tabular}
\end{nota}

% =====================================================================
\section{Ejercicio 1: rendimiento de un fondo (prueba de significancia global)}

\textbf{Contexto.} Como en el ejemplo 1.2 de las notas, un analista explica el rendimiento mensual de
un fondo de inversión, $Y_i$ (\%), con el rendimiento del mercado, $X_{1i}$ (\%), y la tasa de
interés, $X_{2i}$ (\%):
\[
Y_i=\beta_0+\beta_1X_{1i}+\beta_2X_{2i}+\varepsilon_i,\qquad n=10,\quad k=2.
\]
\begin{center}
\begin{tabular}{c*{10}{r}}
\toprule
Mes $i$ & 1 & 2 & 3 & 4 & 5 & 6 & 7 & 8 & 9 & 10\\
\midrule
$X_1$: mercado (\%) & 1.2 & $-0.8$ & 2.0 & 0.5 & $-1.5$ & 1.0 & 2.5 & $-0.3$ & 0.0 & 1.4\\
$X_2$: tasa (\%)    & 7.0 & 7.5 & 6.5 & 7.0 & 8.0 & 7.5 & 6.0 & 7.5 & 7.0 & 6.0\\
$Y$: fondo (\%)     & 1.0 & 0.2 & 1.6 & 1.1 & $-1.2$ & 0.2 & 2.4 & $-0.7$ & 0.9 & 1.5\\
\bottomrule
\end{tabular}
\end{center}

\textbf{Se pide.}
(a) Estimar el modelo por MCO.
(b) Probar al 5\,\% la significancia global $H_0:\beta_1=\beta_2=0$.
(c) Probar al 5\,\% $H_0:\beta_1=0$ con una prueba $F$ y comparar con la prueba $t$.

\subsection*{Solución}

\begin{paso}{Paso 1. Hipótesis del inciso (b)}
\[
H_0:\ \beta_1=\beta_2=0\qquad\text{vs.}\qquad H_1:\ \beta_1\neq0\ \text{o}\ \beta_2\neq0 .
\]
Hay $q=2$ restricciones. Bajo $H_0$ ni el mercado ni la tasa ayudan a explicar el rendimiento del fondo.
\end{paso}

\begin{paso}{Paso 2. Modelo no restringido: medias y desviaciones}
Medias:
\[
\bar X_1=\frac{6.0}{10}=0.6,\qquad \bar X_2=\frac{70.0}{10}=7.0,\qquad \bar Y=\frac{7.0}{10}=0.7 .
\]
Con las desviaciones $x_{1i}=X_{1i}-\bar X_1$, $x_{2i}=X_{2i}-\bar X_2$, $y_i=Y_i-\bar Y$ se arma la tabla:

\begin{center}
\small
\begin{tabular}{c|rrr|rrrrrr}
\toprule
$i$ & $x_1$ & $x_2$ & $y$ & $x_1^2$ & $x_2^2$ & $x_1x_2$ & $x_1y$ & $x_2y$ & $y^2$\\
\midrule
1  & 0.6  & 0.0  & 0.3  & 0.36 & 0.00 & 0.00  & 0.18  & 0.00  & 0.09\\
2  & $-1.4$ & 0.5  & $-0.5$ & 1.96 & 0.25 & $-0.70$ & 0.70  & $-0.25$ & 0.25\\
3  & 1.4  & $-0.5$ & 0.9  & 1.96 & 0.25 & $-0.70$ & 1.26  & $-0.45$ & 0.81\\
4  & $-0.1$ & 0.0  & 0.4  & 0.01 & 0.00 & 0.00  & $-0.04$ & 0.00  & 0.16\\
5  & $-2.1$ & 1.0  & $-1.9$ & 4.41 & 1.00 & $-2.10$ & 3.99  & $-1.90$ & 3.61\\
6  & 0.4  & 0.5  & $-0.5$ & 0.16 & 0.25 & 0.20  & $-0.20$ & $-0.25$ & 0.25\\
7  & 1.9  & $-1.0$ & 1.7  & 3.61 & 1.00 & $-1.90$ & 3.23  & $-1.70$ & 2.89\\
8  & $-0.9$ & 0.5  & $-1.4$ & 0.81 & 0.25 & $-0.45$ & 1.26  & $-0.70$ & 1.96\\
9  & $-0.6$ & 0.0  & 0.2  & 0.36 & 0.00 & 0.00  & $-0.12$ & 0.00  & 0.04\\
10 & 0.8  & $-1.0$ & 0.8  & 0.64 & 1.00 & $-0.80$ & 0.64  & $-0.80$ & 0.64\\
\midrule
$\Sigma$ & 0 & 0 & 0 & 14.28 & 4.00 & $-6.45$ & 10.90 & $-6.05$ & 10.70\\
\bottomrule
\end{tabular}
\end{center}
Es decir: $S_{11}=14.28$, $S_{22}=4.00$, $S_{12}=-6.45$, $S_{1y}=10.90$, $S_{2y}=-6.05$ y
$\SST=S_{yy}=10.70$.
\end{paso}

\begin{paso}{Paso 2 (cont.). Coeficientes de MCO con dos regresoras}
Las ecuaciones normales en desviaciones,
\[
\begin{pmatrix}S_{11}&S_{12}\\ S_{12}&S_{22}\end{pmatrix}
\begin{pmatrix}\hat\beta_1\\ \hat\beta_2\end{pmatrix}
=\begin{pmatrix}S_{1y}\\ S_{2y}\end{pmatrix},
\]
se resuelven con la regla de Cramer. El determinante es
\[
D=S_{11}S_{22}-S_{12}^2=(14.28)(4.00)-(-6.45)^2=57.12-41.6025=15.5175 .
\]
Entonces
\[
\hat\beta_1=\frac{S_{22}S_{1y}-S_{12}S_{2y}}{D}
=\frac{(4.00)(10.90)-(-6.45)(-6.05)}{15.5175}
=\frac{4.5775}{15.5175}=0.2950,
\]
\[
\hat\beta_2=\frac{S_{11}S_{2y}-S_{12}S_{1y}}{D}
=\frac{(14.28)(-6.05)-(-6.45)(10.90)}{15.5175}
=\frac{-16.089}{15.5175}=-1.0368,
\]
\[
\hat\beta_0=\bar Y-\hat\beta_1\bar X_1-\hat\beta_2\bar X_2
=0.7-(0.2950)(0.6)-(-1.0368)(7.0)=7.7808 .
\]
\[
\boxed{\ \hat Y_i=7.7808+0.2950\,X_{1i}-1.0368\,X_{2i}\ }
\]
\textbf{Lectura:} manteniendo fija la tasa, un punto porcentual más de rendimiento de mercado se asocia
con 0.295 puntos más de rendimiento del fondo; manteniendo fijo el mercado, un punto más de tasa se
asocia con 1.037 puntos menos.
\end{paso}

\begin{paso}{Paso 2 (cont.). Suma de cuadrados residual $\SSR_{NR}$}
Atajo con las sumas ya calculadas:
\[
\SSE=\hat\beta_1S_{1y}+\hat\beta_2S_{2y}=(0.29499)(10.90)+(-1.03683)(-6.05)=3.2154+6.2728=9.4882,
\]
\[
\SSR_{NR}=\SST-\SSE=10.70-9.4882=1.2118,\qquad R^2_{NR}=\frac{9.4882}{10.70}=0.8867 .
\]
Comprobación directa con los residuos $\hat e_i=Y_i-\hat Y_i$:
\begin{center}
\small
\footnotesize\setlength{\tabcolsep}{2.5pt}
\begin{tabular}{c*{10}{r}|r}
\toprule
$i$ & 1 & 2 & 3 & 4 & 5 & 6 & 7 & 8 & 9 & 10 & $\Sigma$\\
\midrule
$\hat Y_i$ & 0.8770 & $-0.2314$ & 1.6314 & 0.6705 & $-0.9563$ & 0.2996 & 2.2973 & $-0.0839$ & 0.5230 & 1.9728 & 7.0000\\
$\hat e_i$ & 0.1230 & 0.4314 & $-0.0314$ & 0.4295 & $-0.2437$ & $-0.0996$ & 0.1027 & $-0.6161$ & 0.3770 & $-0.4728$ & 0.0000\\
$\hat e_i^{\,2}$ & 0.0151 & 0.1861 & 0.0010 & 0.1845 & 0.0594 & 0.0099 & 0.0105 & 0.3796 & 0.1421 & 0.2236 & 1.2118\\
\bottomrule
\end{tabular}
\end{center}
Los residuos suman cero (el modelo tiene constante) y $\sum\hat e_i^{\,2}=1.2118$, igual que con el atajo.
\end{paso}

\begin{paso}{Paso 3. Modelo restringido}
Bajo $H_0:\beta_1=\beta_2=0$ el modelo es $Y_i=\beta_0+\varepsilon_i$, cuyo estimador de MCO es
$\hat\beta_0=\bar Y=0.7$. Sus residuos son $Y_i-\bar Y$, así que
\[
\SSR_R=\sum_{i=1}^{10}(Y_i-\bar Y)^2=\SST=10.70 .
\]
\end{paso}

\begin{paso}{Paso 4. Estadístico $F$}
Con $q=2$, $n=10$, $k=2$ y $n-k-1=7$:
\[
F=\frac{(\SSR_R-\SSR_{NR})/q}{\SSR_{NR}/(n-k-1)}
=\frac{(10.70-1.2118)/2}{1.2118/7}
=\frac{4.7441}{0.17311}=27.40 .
\]
Con $R^2$: $\;F=\dfrac{0.8867/2}{(1-0.8867)/7}=\dfrac{0.44337}{0.016179}=27.40$. Coinciden.
\end{paso}

\begin{paso}{Pasos 5 y 6. Decisión y conclusión}
Bajo $H_0$, $F\sim F_{2,7}$ y el valor crítico al 5\,\% es $F_{0.05;\,2,7}=4.737$.
\[
F=27.40>4.737\ \Longrightarrow\ \text{se rechaza }H_0 .
\]
También $27.40>9.547=F_{0.01;\,2,7}$: se rechaza incluso al 1\,\% (valor $p\approx0.0005$).

\textbf{Conclusión:} el rendimiento del mercado y la tasa de interés, considerados conjuntamente,
explican una parte significativa de la variación del rendimiento del fondo.
\end{paso}

\begin{center}
\begin{tikzpicture}
\begin{axis}[width=12cm,height=5cm,domain=0:9,samples=200,xmin=0,xmax=9,ymin=0,ymax=1.1,
  axis lines=left,xlabel={$f$},ylabel={densidad $F_{2,7}$},
  xtick={0,2,4.737,6,8},xticklabels={0,2,4.737,6,8},ytick=\empty,
  title={\small Región de rechazo al 5\,\%: el estadístico observado $F=27.40$ queda muy a la derecha}]
\addplot[name path=f,thick,NavyBlue] {(1+2*x/7)^(-4.5)};
\addplot[draw=none,name path=ax,domain=4.737:9] {0};
\addplot[fill=BrickRed!75,draw=none] fill between[of=f and ax,soft clip={domain=4.737:9}];
\draw[dashed] (axis cs:4.737,0) -- (axis cs:4.737,0.2);
\node[BrickRed,font=\small] at (axis cs:6.6,0.09) {$\alpha=0.05$};
\node[font=\small,anchor=east] at (axis cs:8.9,0.2) {$F=27.40\ \longrightarrow$};
\end{axis}
\end{tikzpicture}
\end{center}

\begin{paso}{Inciso (c). Una sola restricción: $H_0:\beta_1=0$}
\textbf{Hipótesis:} $H_0:\beta_1=0$ vs.\ $H_1:\beta_1\neq0$; ahora $q=1$.

\textbf{Modelo restringido:} quitar $X_1$ y estimar $Y_i=\gamma_0+\gamma_2X_{2i}+\varepsilon_i$
(regresión simple):
\[
\hat\gamma_2=\frac{S_{2y}}{S_{22}}=\frac{-6.05}{4.00}=-1.5125,\qquad
\hat\gamma_0=\bar Y-\hat\gamma_2\bar X_2=0.7+1.5125(7.0)=11.2875,
\]
\[
\SSR_R=S_{yy}-\hat\gamma_2S_{2y}=10.70-(-1.5125)(-6.05)=10.70-9.1506=1.5494 .
\]
\textbf{Estadístico:}
\[
F=\frac{(1.5494-1.2118)/1}{1.2118/7}=\frac{0.3376}{0.17311}=1.950 .
\]
\textbf{Decisión:} $F_{0.05;\,1,7}=5.591$ y $1.950<5.591$: no se rechaza $H_0$ (valor $p\approx0.21$).

\textbf{Relación con la prueba $t$.} Con $\hat\sigma^2=\SSR_{NR}/(n-k-1)=0.17311$,
\[
\operatorname{se}(\hat\beta_1)=\sqrt{\frac{\hat\sigma^2\,S_{22}}{D}}
=\sqrt{\frac{(0.17311)(4.00)}{15.5175}}=0.2112,\qquad
t=\frac{0.2950}{0.2112}=1.396,
\]
y $t^2=(1.396)^2=1.950=F$. Lo mismo ocurre con los valores críticos:
$t_{0.025;7}^2\approx5.591=F_{0.05;1,7}$ (el valor crítico $t$ redondeado es $2.365$). Con una sola restricción, la prueba $F$ y la prueba $t$
bilateral son la misma prueba.
\end{paso}

\begin{nota}{Lección del ejercicio 1}
El mercado no es significativo \emph{individualmente} ($F=1.95$), pero las dos regresoras sí lo son
\emph{conjuntamente} ($F=27.40$). No hay contradicción: el mercado y la tasa están fuertemente
correlacionados en la muestra (correlación $-0.85$; observe que $S_{12}=-6.45$), de modo que los datos
apenas permiten separar el efecto de cada una, aunque dejan claro que el bloque explica al fondo. Por
eso una prueba conjunta no se reemplaza por varias pruebas $t$ por separado.
\end{nota}

% =====================================================================
\section{Ejercicio 2: Size, Value y Momentum (prueba de exclusión de un bloque)}

\textbf{Contexto.} Como en el ejemplo 3.2 de las notas, se explica el rendimiento anual en exceso de
una acción, $R_i$ (\%), con tres características: tamaño $S_i$ (\emph{Size}, escala 1--5), valor $V_i$
(\emph{Value}, escala 0--2) y momento $M_i$ (\emph{Momentum}, escala $-2$ a 3):
\[
R_i=\beta_0+\beta_1S_i+\beta_2V_i+\beta_3M_i+\varepsilon_i,\qquad n=12,\quad k=3 .
\]
\begin{center}
\begin{tabular}{c*{12}{r}}
\toprule
Acción $i$ & 1 & 2 & 3 & 4 & 5 & 6 & 7 & 8 & 9 & 10 & 11 & 12\\
\midrule
$S$ & 2 & 4 & 1 & 3 & 5 & 2 & 4 & 3 & 1 & 5 & 3 & 4\\
$V$ & 1 & 0 & 2 & 1 & 0 & 2 & 1 & 0 & 2 & 1 & 1 & 0\\
$M$ & 3 & 1 & 2 & $-1$ & 0 & 1 & 2 & $-2$ & 0 & 1 & $-1$ & 2\\
$R$ (\%) & 3.0 & 1.5 & 4.0 & 1.0 & $-0.5$ & 2.5 & 1.0 & 0.5 & 2.5 & 0.0 & 2.5 & 1.0\\
\bottomrule
\end{tabular}
\end{center}

\textbf{Se pide.} Probar al 5\,\% si Value y Momentum, como bloque, aportan información sobre el
rendimiento una vez que se controla por tamaño: $H_0:\beta_2=\beta_3=0$.

\subsection*{Solución}

\begin{paso}{Paso 1. Hipótesis}
\[
H_0:\ \beta_2=\beta_3=0\qquad\text{vs.}\qquad H_1:\ \beta_2\neq0\ \text{o}\ \beta_3\neq0,
\qquad q=2 .
\]
\end{paso}

\begin{paso}{Paso 2. Modelo no restringido en forma matricial}
La matriz de diseño $X$ tiene una columna de unos y las columnas $S$, $V$, $M$. Sus productos cruzados
son sumas que se obtienen directamente de la tabla de datos:
\[
X'X=\begin{pmatrix}
n & \sum S & \sum V & \sum M\\
\sum S & \sum S^2 & \sum SV & \sum SM\\
\sum V & \sum SV & \sum V^2 & \sum VM\\
\sum M & \sum SM & \sum VM & \sum M^2
\end{pmatrix}
=\begin{pmatrix}
12 & 37 & 11 & 8\\
37 & 135 & 25 & 23\\
11 & 25 & 17 & 10\\
8 & 23 & 10 & 30
\end{pmatrix},
\]
\[
X'\mathbf y=\begin{pmatrix}\sum R\\ \sum SR\\ \sum VR\\ \sum MR\end{pmatrix}
=\begin{pmatrix}19.0\\ 41.0\\ 25.5\\ 20.5\end{pmatrix}.
\]
Por ejemplo, $\sum SV=2(1)+4(0)+1(2)+\cdots+4(0)=25$ y $\sum MR=3(3.0)+1(1.5)+\cdots+2(1.0)=20.5$.
El determinante es $\det(X'X)=18\,225\neq0$ (no hay colinealidad perfecta) y la inversa es
\[
(X'X)^{-1}=\begin{pmatrix}
\phantom{-}2.14579 & -0.45663 & -0.72927 & \phantom{-}0.02096\\
-0.45663 & \phantom{-}0.10765 & \phantom{-}0.14189 & -0.00807\\
-0.72927 & \phantom{-}0.14189 & \phantom{-}0.33789 & -0.02694\\
\phantom{-}0.02096 & -0.00807 & -0.02694 & \phantom{-}0.04291
\end{pmatrix}.
\]
(Verifique, por ejemplo, que la primera fila de $X'X$ por la primera columna de la inversa da 1.)
Entonces $\hat{\boldsymbol\beta}=(X'X)^{-1}X'\mathbf y$; la primera pendiente es
\[
\hat\beta_1=(-0.45663)(19.0)+(0.10765)(41.0)+(0.14189)(25.5)+(-0.00807)(20.5)=-0.8091,
\]
y del mismo modo se obtienen las demás:
\[
\boxed{\ \hat R_i=3.8816-0.8091\,S_i+0.0253\,V_i+0.2602\,M_i\ }
\]
\end{paso}

\begin{paso}{Paso 2 (cont.). Suma de cuadrados residual $\SSR_{NR}$}
Atajo matricial: $\SSR=\mathbf y'\mathbf y-\hat{\boldsymbol\beta}'X'\mathbf y$, con
$\mathbf y'\mathbf y=\sum R_i^2=49.50$ y
\[
\hat{\boldsymbol\beta}'X'\mathbf y=(3.88156)(19.0)+(-0.80914)(41.0)+(0.02527)(25.5)+(0.26016)(20.5)=46.5528,
\]
\[
\SSR_{NR}=49.50-46.5528=2.9472 .
\]
Comprobación con los residuos:
\begin{center}
\small
\footnotesize\setlength{\tabcolsep}{2pt}
\begin{tabular}{c*{12}{r}|r}
\toprule
$i$ & 1 & 2 & 3 & 4 & 5 & 6 & 7 & 8 & 9 & 10 & 11 & 12 & $\Sigma$\\
\midrule
$\hat R_i$ & 3.069 & 0.905 & 3.643 & 1.219 & $-0.164$ & 2.574 & 1.191 & 0.934 & 3.123 & 0.121 & 1.219 & 1.165 & 19.000\\
$\hat e_i$ & $-0.069$ & 0.595 & 0.357 & $-0.219$ & $-0.336$ & $-0.074$ & $-0.191$ & $-0.434$ & $-0.623$ & $-0.121$ & 1.281 & $-0.165$ & 0.000\\
$\hat e_i^{\,2}$ & 0.005 & 0.354 & 0.127 & 0.048 & 0.113 & 0.005 & 0.036 & 0.188 & 0.388 & 0.015 & 1.640 & 0.027 & 2.947\\
\bottomrule
\end{tabular}
\end{center}
\end{paso}

\begin{paso}{Paso 3. Modelo restringido: regresión simple sobre $S$}
Imponer $\beta_2=\beta_3=0$ equivale a estimar $R_i=\gamma_0+\gamma_1S_i+\varepsilon_i$. Con las
fórmulas de cálculo (que usan las mismas sumas de $X'X$ y $X'\mathbf y$):
\[
S_{SS}=\sum S^2-\frac{(\sum S)^2}{n}=135-\frac{37^2}{12}=135-114.0833=20.9167,
\]
\[
S_{SR}=\sum SR-\frac{(\sum S)(\sum R)}{n}=41.0-\frac{(37)(19.0)}{12}=41.0-58.5833=-17.5833,
\]
\[
S_{RR}=\sum R^2-\frac{(\sum R)^2}{n}=49.50-\frac{19.0^2}{12}=49.50-30.0833=19.4167=\SST .
\]
Coeficientes:
\[
\hat\gamma_1=\frac{S_{SR}}{S_{SS}}=\frac{-17.5833}{20.9167}=-0.8406,
\]
\[
\hat\gamma_0=\bar R-\hat\gamma_1\bar S=\frac{19}{12}+0.8406\Big(\frac{37}{12}\Big)=1.5833+2.5920=4.1753 .
\]
Suma de cuadrados residual:
\[
\SSR_R=S_{RR}-\hat\gamma_1S_{SR}=19.4167-(-0.8406)(-17.5833)=19.4167-14.7812=4.6355 .
\]
\end{paso}

\begin{paso}{Paso 4. Estadístico $F$}
Con $q=2$, $n=12$, $k=3$ y $n-k-1=8$:
\[
F=\frac{(\SSR_R-\SSR_{NR})/q}{\SSR_{NR}/(n-k-1)}
=\frac{(4.6355-2.9472)/2}{2.9472/8}
=\frac{0.84415}{0.36840}=2.291 .
\]
Con $R^2$: $R^2_{NR}=1-2.9472/19.4167=0.8482$ y $R^2_R=1-4.6355/19.4167=0.7613$, de modo que
\[
F=\frac{(0.8482-0.7613)/2}{(1-0.8482)/8}=\frac{0.043476}{0.018974}=2.291 .
\]
\end{paso}

\begin{paso}{Pasos 5 y 6. Decisión y conclusión}
Bajo $H_0$, $F\sim F_{2,8}$ y $F_{0.05;\,2,8}=4.459$.
\[
F=2.291<4.459\ \Longrightarrow\ \text{no se rechaza }H_0\quad(\text{valor }p\approx0.16).
\]
Ni siquiera al 10\,\% se rechaza, pues $2.291<3.113=F_{0.10;\,2,8}$.

\textbf{Conclusión:} una vez que se controla por tamaño, Value y Momentum, como bloque, no aportan
información estadísticamente significativa sobre el rendimiento en esta muestra. Quitarlas eleva la
suma de cuadrados residual de 2.947 a 4.636, pero ese aumento es compatible con el ruido residual.
\end{paso}

\begin{center}
\begin{tikzpicture}
\begin{axis}[width=12cm,height=5cm,domain=0:9,samples=200,xmin=0,xmax=9,ymin=0,ymax=1.1,
  axis lines=left,xlabel={$f$},ylabel={densidad $F_{2,8}$},
  xtick={0,2.291,4.459,6,8},xticklabels={0,2.291,4.459,6,8},ytick=\empty,
  title={\small El estadístico observado $F=2.291$ cae en la región de no rechazo}]
\addplot[name path=f,thick,NavyBlue] {(1+x/4)^(-5)};
\addplot[draw=none,name path=ax,domain=4.459:9] {0};
\addplot[fill=BrickRed!75,draw=none] fill between[of=f and ax,soft clip={domain=4.459:9}];
\draw[dashed] (axis cs:4.459,0) -- (axis cs:4.459,0.2);
\draw[very thick,OliveGreen] (axis cs:2.291,0) -- (axis cs:2.291,0.17);
\node[OliveGreen,font=\small,anchor=south] at (axis cs:2.291,0.17) {$F$ observado};
\node[BrickRed,font=\small] at (axis cs:6.4,0.09) {$\alpha=0.05$};
\end{axis}
\end{tikzpicture}
\end{center}

\begin{nota}{Lección del ejercicio 2}
Que el bloque $\{V,M\}$ no sea significativo no significa que el modelo completo sea inútil. La
prueba global del mismo modelo, $H_0:\beta_1=\beta_2=\beta_3=0$, da
\[
F=\frac{\SSE/k}{\SSR/(n-k-1)}=\frac{(19.4167-2.9472)/3}{2.9472/8}=14.90>4.066=F_{0.05;\,3,8},
\]
así que el modelo sí explica el rendimiento; lo que explica proviene esencialmente del tamaño. Cada
prueba $F$ responde a la pregunta que fija su $H_0$: hay que leer siempre \emph{qué} restricciones se
están probando.
\end{nota}

\section{Ejercicio 3: prueba $t$ individual e intervalo de confianza}

\textbf{Contexto.} Retomamos la regresión del ejercicio 1, que explica el rendimiento del fondo
con el rendimiento del mercado y la tasa de interés. La pendiente estimada de la tasa y su error
estándar fueron $\hat\beta_2=-1.0368$ y
$\operatorname{se}(\hat\beta_2)=0.3991$, respectivamente.
Con $n=10$ observaciones y $k=2$ regresoras, los grados de libertad residuales son
$\nu=n-k-1=7$.

\textbf{Se pide.} (a) Probar al 5\,\% $H_0:\beta_2=0$ contra $H_1:\beta_2\neq0$ con una prueba
$t$. (b) Construir el intervalo de confianza del 95\,\% para $\beta_2$. (c) Representar la regla
de decisión en la distribución $t_7$.

\subsection*{Solución}

\begin{paso}{Paso 1. Hipótesis y estadístico}
La prueba es bilateral porque la alternativa permite que la pendiente sea positiva o negativa:
\[
H_0:\beta_2=0,\qquad H_1:\beta_2\neq0,\qquad \alpha=0.05.
\]
El estadístico compara la estimación con el valor indicado por $H_0$, medido en errores estándar:
\[
t=\frac{\hat\beta_2-0}{\operatorname{se}(\hat\beta_2)}
=\frac{-1.0368}{0.3991}=-2.598,\qquad \nu=7.
\]
Bajo $H_0$ y los supuestos clásicos, $T\sim t_7$.
\end{paso}

\begin{paso}{Paso 2. Valor crítico y decisión}
Al ser bilateral, se asigna $\alpha/2=0.025$ a cada cola. El valor crítico es
$t_{0.025;\,7}=2.365$:
\[
|t|=2.598>2.365\ \Longrightarrow\ \text{se rechaza }H_0.
\]
La muestra aporta evidencia al nivel del 5\,\% de que la pendiente de la tasa de interés difiere
de cero. Su signo negativo indica una asociación inversa con el rendimiento del fondo, manteniendo
fijo el rendimiento del mercado.
\end{paso}

\begin{paso}{Paso 3. Intervalo de confianza del 95\,\%}
El intervalo es la estimación más y menos el valor crítico multiplicado por su error estándar:
\[
\hat\beta_2\pm t_{0.025;\,7}\operatorname{se}(\hat\beta_2)
=-1.0368\pm(2.365)(0.3991)
=-1.0368\pm0.9439,
\]
\[
\boxed{\ \beta_2\in[-1.9807,\;-0.0929]\quad\text{(intervalo de confianza del 95\,\%)}\ }.
\]
El intervalo no contiene cero, en concordancia con la decisión de rechazar $H_0$.
\end{paso}

\begin{center}
\begin{tikzpicture}
\begin{axis}[
  width=12cm,height=5.5cm,domain=-4:4,samples=250,
  xmin=-4,xmax=4,ymin=0,ymax=0.42,
  axis lines=left,xlabel={$t$},ylabel={densidad $t_7$},
  xtick={-2.598,-2.365,0,2.365},
  xticklabels={$-2.598$,$-2.365$,0,2.365},ytick=\empty,
  title={\small Prueba bilateral al 5\,\%: el estadístico observado cae en la región de rechazo}]
\addplot[name path=student,thick,NavyBlue]
  {0.389108*(1+x^2/7)^(-4)};
\addplot[draw=none,name path=axisleft,domain=-4:-2.365] {0};
\addplot[fill=BrickRed!75,draw=none]
  fill between[of=student and axisleft,soft clip={domain=-4:-2.365}];
\addplot[draw=none,name path=axisright,domain=2.365:4] {0};
\addplot[fill=BrickRed!75,draw=none]
  fill between[of=student and axisright,soft clip={domain=2.365:4}];
\draw[dashed,BrickRed] (axis cs:-2.365,0) -- (axis cs:-2.365,0.12);
\draw[dashed,BrickRed] (axis cs:2.365,0) -- (axis cs:2.365,0.12);
\draw[very thick,OliveGreen] (axis cs:-2.598,0) -- (axis cs:-2.598,0.32);
\node[OliveGreen,font=\small,anchor=south] at (axis cs:-2.598,0.32) {$t=-2.598$};
\node[BrickRed,font=\small] at (axis cs:3.1,0.10) {$\alpha/2=0.025$};
\node[BrickRed,font=\small] at (axis cs:-3.1,0.10) {$\alpha/2=0.025$};
\end{axis}
\end{tikzpicture}
\end{center}

\begin{nota}{Comprobación con la prueba $F$}
Para una sola restricción bilateral, la prueba $F$ equivalente da
$F_{1,7}=t^2\approx6.748$. Como $6.748>F_{0.05;\,1,7}=5.591$, también se rechaza $H_0$.
Esto ilustra que $F=t^2$ conserva la decisión de la prueba $t$ bilateral, aunque no conserva el
signo del efecto.
\end{nota}

\section{Resumen de los tres ejercicios}

\begin{center}
\small
\begin{tabular}{llcccc}
\toprule
Ejercicio & Hipótesis & Distribución & Estadístico & Valor crítico (5\,\%) & Decisión\\
\midrule
E1(b) & $\beta_1=\beta_2=0$ & $F_{2,7}$ & 27.40 & 4.737 & Rechazar\\
E1(c) & $\beta_1=0$ & $F_{1,7}$ & 1.950 & 5.591 & No rechazar\\
E2 & $\beta_2=\beta_3=0$ & $F_{2,8}$ & 2.291 & 4.459 & No rechazar\\
E2 global & $\beta_1=\beta_2=\beta_3=0$ & $F_{3,8}$ & 14.90 & 4.066 & Rechazar\\
E3 & $\beta_2=0$ & $t_7$ bilateral & $-2.598$ & $\pm2.365$ & Rechazar\\
\bottomrule
\end{tabular}
\end{center}

\textit{Nota:} para una prueba bilateral, la regla $|t|>t_{\alpha/2;\nu}$ equivale a
$t^2>F_{\alpha;1,\nu}$.

\begin{nota}{Errores frecuentes}
\begin{itemize}[leftmargin=1.2em,itemsep=1pt]
\item Usar la suma de cuadrados explicada en lugar de la \emph{residual} en la fórmula de $F$.
\item Tomar $k$ del modelo restringido: los grados de libertad del denominador son $n-k-1$ del modelo
\emph{no restringido}.
\item Contar mal $q$: es el número de signos ``$=$'' en $H_0$, no el número de regresoras del modelo.
\item Usar la forma con $R^2$ cuando los modelos tienen distinta variable dependiente (por ejemplo
$Y$ en uno y $\ln Y$ en el otro); en ese caso sólo vale la forma con $\SSR$.
\item Concluir que ``cada coeficiente es significativo'' cuando se rechaza una $H_0$ conjunta: el
rechazo dice que \emph{al menos una} restricción falla.
\end{itemize}
\end{nota}

\end{document}
